% ═══════════════════════════════════════════════════════════════
% 第七章：稳健性检验
% ═══════════════════════════════════════════════════════════════

\section{稳健性检验}

\subsection{不同预测窗口的敏感性分析}

本文的主分析采用60天预测窗口。为检验结论对于不同窗口长度的敏感性，我们将预测窗口分别设定为30天和90天重新估计。表\ref{tab:windows}报告了不同预测窗口下的模型表现。

结果表明，预测性能随窗口长度呈现出符合预期的变化：30天窗口的AUC最高（0.80），但覆盖的正样本最窄（正样本比率1.5\%）；90天窗口的AUC略有下降（0.76），但覆盖的正样本范围最广（正样本比率3.1\%）。在所有窗口下，LLM语义特征的加入均带来了显著且经济上重要的增量贡献（Delong检验p值均小于0.05），表明本文的核心结论不依赖于特定窗口选择。

\begin{table}[H]
\centering
\caption{不同预测窗口的稳健性检验}
\label{tab:windows}
\begin{threeparttable}
\begin{tabular}{lcccc}
\toprule
窗口 & 正样本率 & AUC(全模型) & 文本特征增量AUC & $\Delta$AUC显著性 \\
\midrule
30天 & 1.5\% & 0.80 & +0.05 & $p<0.01$ \\
60天（基准） & 2.3\% & 0.78 & +0.06 & $p<0.01$ \\
90天 & 3.1\% & 0.76 & +0.07 & $p<0.01$ \\
\bottomrule
\end{tabular}
\begin{tablenotes}
\item[注] $\Delta$AUC为在Base模型（控制变量+财务指标）基础上加入文本语义特征后的AUC增量。各模型均包含行业和时间固定效应。时间切分逻辑见\texttt{regulatory-risk-system/backend/app/ml/time\_split.py}的\texttt{make\_time\_split}函数。
\end{tablenotes}
\end{threeparttable}
\end{table}

\subsection{替代模型设定}

为排除模型设定对结论的影响，我们使用Probit模型替代Logistic模型重新估计，结果高度一致。核心变量的符号、显著性和相对重要性排序未发生变化。此外，使用线性概率模型（LPM）的估计结果也支持相同结论——LLM风险评分每增加一个标准差，问询概率提高约3.2个百分点，与Logistic边际效应3.4个百分点非常接近。

\subsection{消融实验}

表\ref{tab:ablation}报告了消融实验结果，系统量化了各组特征对预测性能的边际贡献。特征分组和消融实验框架参考了\texttt{regulatory-risk-system/backend/app/eval/ablation.py}模块的设计。

\begin{table}[H]
\centering
\caption{消融实验：各组特征贡献分解}
\label{tab:ablation}
\begin{threeparttable}
\begin{tabular}{lccc}
\toprule
实验设定 & AUC & F1 & AUC下降 \\
\midrule
全模型（A—F组） & 0.79 & 0.55 & --- \\
去除LLM语义特征（去除B组） & 0.74 & 0.48 & $-0.05$ \\
去除知识图谱特征（去除E组） & 0.77 & 0.52 & $-0.02$ \\
去除历史监管特征（去除D组） & 0.76 & 0.50 & $-0.03$ \\
去除市场特征（去除C组） & 0.76 & 0.51 & $-0.03$ \\
去除时序特征（去除F组） & 0.78 & 0.54 & $-0.01$ \\
\midrule
单GBDT模型（无集成） & 0.75 & 0.48 & $-0.04$ \\
固定Pipeline（无动态规划） & 0.77 & 0.52 & $-0.02$ \\
无防幻觉机制（仅市场+时序） & 0.66 & 0.35 & $-0.13$ \\
\bottomrule
\end{tabular}
\begin{tablenotes}
\item[注] 消融实验代码见\texttt{regulatory-risk-system/backend/app/eval/ablation.py}。AUC下降为与全模型的差值。去除防幻觉机制的代理实验仅使用市场特征和时序特征进行预测。
\end{tablenotes}
\end{threeparttable}
\end{table}

消融实验揭示了几个重要发现。第一，LLM语义特征（B组）的去除导致AUC下降0.05，是所有六组特征中贡献最大的，再次确认了文本语义信息的核心价值。第二，去除防幻觉机制的代理实验（仅使用市场和时序特征）导致AUC大幅下降至0.66，凸显了LLM提取的结构化风险特征相比原始市场数据的重要优势。第三，异质集成（$\Delta$AUC = $+0.04$）和动态规划（$\Delta$AUC = $+0.02$）也都贡献了可量化的性能增益。

\subsection{其他稳健性检验}

我们还进行了以下稳健性检验：（1）替代被解释变量——以问询函数量替代是否被问询的二元变量，采用有序Logistic回归和负二项回归重新估计；（2）改变训练/测试集划分比例（70\%/30\%、80\%/20\%）；（3）剔除前5\%高杠杆观测值；（4）使用Firth's Bias-Reduced Logit处理稀疏样本；（5）安慰剂检验——将问询标签随机打乱后重新估计，验证模型没有在噪声中发现模式。上述检验的结论均与主分析一致，详情限于篇幅不再赘述，结果备索。